\documentclass[11pt]{article}
\usepackage[margin=1in]{geometry}
\usepackage{amsmath,amssymb,amsthm}
\usepackage[hidelinks]{hyperref}

\newtheorem{theorem}{Theorem}
\newtheorem{lemma}{Lemma}
\newtheorem{corollary}{Corollary}

\newcommand{\N}{\mathbb N}
\newcommand{\R}{\mathbb R}
\newcommand{\linf}{\ell_\infty}
\newcommand{\co}{\operatorname{co}}

\title{A Counterexample to Sequential-ASQ Renormability for All Spaces Containing \(c_0\)}
\author{Run \texttt{fa\_banach\_001}, agent lane 17}
\date{2026-06-20}

\begin{document}
\maketitle

\section*{Status}

This packet gives a candidate ZFC counterexample, likely valid, to the
renorming question asked in Aviles--Martinez-Cervantes--Rueda Zoca
\cite[p.~2]{AMRZ}.  The source paper asks whether every Banach space containing
\(c_0\) can be equivalently renormed to be sequentially ASQ.  We show that the
answer is no: \(\ell_\infty\) contains \(c_0\), but no equivalent real norm on
\(\ell_\infty\) is sequentially ASQ.

\section*{Proof Intuition}

The obstruction is that a sequentially ASQ witness carries hidden
complementability.  Let \((x_n)\subset S_X\) witness sequential ASQ, and choose
norming functionals \(f_n(x_n)=1\).  Since
\[
        \|x_n\pm t x\|\longrightarrow 1
        \qquad (0<t\|x\|<1),
\]
the inequalities obtained by applying \(f_n\) force \(f_n(x)\to 0\) for each
fixed \(x\in X\).  After passing to a diagonal subsequence, the matrix
\((f_i(x_j))\) is a small perturbation of the identity on \(c_0\).  The map
\(x\mapsto (f_i(x))\) therefore gives, after inverting this matrix, a bounded
projection from \(X\) onto the \(c_0\)-copy generated by the chosen subsequence.

\section*{The Obstruction Lemma}

\begin{lemma}\label{lem:pointwise}
Let \(X\) be a real sequentially ASQ Banach space, witnessed by
\((x_n)\subset S_X\).  For each \(n\), choose \(f_n\in S_{X^*}\) with
\(f_n(x_n)=1\).  Then \(f_n(x)\to 0\) for every \(x\in X\).  Also
\((x_n)\) is weakly null.
\end{lemma}

\begin{proof}
Fix \(x\in X\).  If \(x=0\) there is nothing to prove.  Choose
\(t>0\) with \(t\|x\|<1\).  By sequential ASQ, applied to \(t x\) and to
\(-t x\),
\[
        \|x_n+t x\|\to 1,
        \qquad
        \|x_n-t x\|\to 1.
\]
Since \(f_n(x_n)=1\) and \(\|f_n\|=1\),
\[
        1+t f_n(x) \leq \|x_n+t x\|,
        \qquad
        1-t f_n(x) \leq \|x_n-t x\|.
\]
Letting \(n\to\infty\) gives \(t f_n(x)\to 0\), hence \(f_n(x)\to 0\).

Now let \(g\in X^*\).  We prove \(g(x_n)\to0\).  By homogeneity assume
\(\|g\|=1\).  If not, for some \(\delta>0\) there is a subsequence along which
\(|g(x_n)|\geq\delta\).  Passing to a further subsequence and replacing
\(g\) by \(-g\) if necessary, assume \(g(x_n)\geq\delta\).  Pick \(u\in S_X\)
with \(g(u)>1-\eta\), where \(\eta>0\) is small.  For any \(\lambda>1\),
sequential ASQ gives \(\|\lambda u+x_n\|\to\lambda\), while along the chosen
subsequence
\[
        \|\lambda u+x_n\|
        \geq g(\lambda u+x_n)
        \geq \lambda(1-\eta)+\delta.
\]
Taking \(\eta<\delta/(2\lambda)\) contradicts convergence to \(\lambda\).
Thus \(g(x_n)\to0\), as claimed.
\end{proof}

\begin{lemma}\label{lem:complemented}
Every real sequentially ASQ Banach space contains a complemented subspace
isomorphic to \(c_0\).
\end{lemma}

\begin{proof}
Let \((x_n)\subset S_X\) witness sequential ASQ, and choose
\(f_n\in S_{X^*}\) with \(f_n(x_n)=1\).  By Lemma \ref{lem:pointwise},
\(f_n(x)\to0\) for each fixed \(x\), and \(x_n\to0\) weakly.

We pass to a subsequence, still denoted \((u_k)\), with associated
functionals \((g_k)\), satisfying two standard diagonal requirements.
First, using the finite-dimensional selection lemma for sequentially ASQ
sequences from \cite[Lemma 3.2]{AMRZ}, we make \((u_k)\) \(2\)-equivalent
from above to the \(c_0\)-basis:
\[
        \left\|\sum_{k=1}^m a_k u_k\right\|
        \leq 2\max_{1\leq k\leq m}|a_k|
        \qquad(m\in\N).
\]
Second, using \(g_n(x)\to0\) for fixed \(x\) and weak nullity of \((x_n)\), we
diagonalize further so that, for all \(i\neq j\),
\[
        |g_i(u_j)| < 2^{-i-j-3}.
\]
Indeed, at stage \(k\) only finitely many previous vectors and previous
functionals are involved, so the two pointwise convergence statements allow
choosing the next index large enough.

Let \(Y=\overline{\operatorname{span}}\{u_k:k\in\N\}\).  Define
\[
        T:X\longrightarrow c_0,\qquad T x=(g_k(x))_{k=1}^\infty.
\]
Lemma \ref{lem:pointwise} gives \(T x\in c_0\) for every \(x\), and
\(\|T\|\leq1\).  Also define \(U:c_0\to Y\) by
\[
        U(a)=\sum_{k=1}^\infty a_k u_k .
\]
The upper \(c_0\)-estimate above makes \(U\) bounded.

Consider \(A=TU:c_0\to c_0\).  Its matrix is \(A_{ij}=g_i(u_j)\).  The
diagonal entries are \(1\), and the off-diagonal estimate gives
\[
        \|A-I_{c_0}\|
        \leq \sup_i \sum_{j\neq i}|g_i(u_j)|
        <\frac12.
\]
Hence \(A\) is invertible on \(c_0\).  It follows that \(U\) is an isomorphic
embedding of \(c_0\) into \(X\), and
\[
        P=U A^{-1}T:X\longrightarrow Y
\]
is a bounded projection, since \(P U=U A^{-1}TU=U\).  Therefore \(Y\) is a
complemented copy of \(c_0\) in \(X\).
\end{proof}

\section*{Counterexample}

\begin{theorem}
The real Banach space \(\ell_\infty\) contains an isomorphic copy of \(c_0\)
but does not admit any equivalent sequentially ASQ norm.
\end{theorem}

\begin{proof}
The canonical inclusion \(c_0\subset\ell_\infty\) shows that
\(\ell_\infty\) contains \(c_0\).  Suppose that an equivalent norm
\(|\cdot|\) on \(\ell_\infty\) were sequentially ASQ.  By Lemma
\ref{lem:complemented}, the renormed space \((\ell_\infty,|\cdot|)\) would
contain a complemented subspace isomorphic to \(c_0\).  Complementability of a
closed subspace is preserved under equivalent renorming of the ambient Banach
space, so the usual \(\ell_\infty\) would contain a complemented copy of
\(c_0\).

This contradicts Phillips' classical theorem, in the standard form that every
bounded operator \(\ell_\infty\to c_0\) is weakly compact \cite{Phillips}.
Indeed, a complemented copy \(Y\subset \ell_\infty\) with \(Y\simeq c_0\)
would give an onto operator \(\ell_\infty\to c_0\), obtained by composing the
projection onto \(Y\) with an isomorphism \(Y\to c_0\).  By the open mapping
theorem, an onto weakly compact operator would make the unit ball of \(c_0\)
relatively weakly compact, forcing \(c_0\) to be reflexive.  This is impossible.
Hence no equivalent norm on \(\ell_\infty\) is sequentially ASQ.
\end{proof}

\section*{Consequence for the Source Question}

The source paper's question ``whether every Banach space containing \(c_0\)
can be renormed to be sequentially ASQ'' has a negative ZFC answer.  The
counterexample is \(\ell_\infty\).  This does not address the separate
set-theoretic question in \cite{AMRZ} about removing selective ultrafilters
from their fourth-dual theorem for already sequentially ASQ spaces.

\section*{Novelty and Search Bounds}

The run indexes were searched for \texttt{2110.14313}, \texttt{sequentially
ASQ}, \texttt{sequentially almost square}, \texttt{renormed to be
sequentially ASQ}, \texttt{complemented c0}, and \(\ell_\infty\).  The only
nearby packet records the known nonsequential ASQ renorming theorem for spaces
containing \(c_0\).  Web searches on 2026-06-20 for exact phrases around the
source question and the obstruction above found the source paper but no later
answer.

\begin{thebibliography}{9}

\bibitem{AMRZ}
A. Aviles, G. Martinez-Cervantes, and A. Rueda Zoca,
\emph{Banach spaces containing \(c_0\) and elements in the fourth dual},
arXiv:2110.14313, 2021.

\bibitem{Phillips}
R. S. Phillips,
\emph{On linear transformations},
Transactions of the American Mathematical Society \textbf{48} (1940),
516--541.

\end{thebibliography}

\end{document}
